%% ============================================================
%%  Chapter 2: Background and Related Work
%% ============================================================
\chapter{Background and Related Work}
\label{chap:background}

% GUIDANCE: 8--12 pages. Organise as a structured review, NOT a
% list of paper summaries. Each section should end with a short
% paragraph positioning your work relative to what was reviewed.
% Write this chapter early -- it is largely independent of results.


% ── 2.1 ─────────────────────────────────────────────────────────
\section{Multi-Robot Task Allocation}
\label{sec:bg_mrta}

% Cover the taxonomy of MRTA: ST-SR-IA, MT-SR-IA, etc.
% Ground the reader in why task allocation is NP-hard in general
% and what approximate methods are used in practice.

\todo{Introduce MRTA taxonomy. Cover: single-task vs multi-task
robots, single-robot vs multi-robot tasks, instantaneous vs
time-extended assignment. Cite the Gerkey and Mataric taxonomy.}


% ── 2.1.1 ──────────────────────────────────────────────────────
\subsection{Centralised Approaches}
\label{subsec:bg_centralised}

% Cover: Hungarian algorithm, ILP formulations, SGA.
% Identify the scalability and single-point-of-failure problems
% that motivate decentralised alternatives.

\todo{Describe centralised MRTA: optimality, computational cost,
fragility under communication failures. SGA as the relevant
benchmark~\citep{kha2025enhancing}.}


% ── 2.1.2 ──────────────────────────────────────────────────────
\subsection{Decentralised Auction-Based Approaches}
\label{subsec:bg_decentralised}

% Core content: CBBA (Choi et al. 2009), its guarantees,
% its limitations on non-fully-connected graphs.
% Then: GCBBA (Kha et al. 2025) -- the three improvements:
%   (1) RPT metric is DMG, (2) global consensus restores SGA
%       convergence, (3) decentralised early convergence detection.

\todo{Describe CBBA~\citep{choi2009consensus}: bundle building,
conflict resolution, 50\% optimality under DMG. Then explain
GCBBA~\citep{kha2025enhancing}: why CBBA fails to converge to SGA
on disconnected graphs, and how GCBBA fixes this via global
consensus. This is the foundational algorithm of your thesis.}


% ── 2.1.3 ──────────────────────────────────────────────────────
\subsection{Other Decentralised Methods}
\label{subsec:bg_other_decentralised}

% Brief coverage: market-based methods, DSA, MaxSum for DCOPs.
% The point is to show you know the landscape, not to review
% every paper.

\todo{Brief survey of other decentralised approaches.
Cite representative work. Conclude with why GCBBA is the
most relevant starting point for your problem.}


% ── 2.2 ─────────────────────────────────────────────────────────
\section{Multi-Agent Path Finding}
\label{sec:bg_mapf}

% Cover: the MAPF problem definition, optimality, completeness.
% Vertex and edge conflict definitions (critical for your work).

\todo{Define MAPF formally. Cover completeness and optimality
criteria. Introduce the distinction between offline (plan-then-
execute) and online/lifelong MAPF.}


% ── 2.2.1 ──────────────────────────────────────────────────────
\subsection{Classical MAPF Solvers}
\label{subsec:bg_mapf_classical}

% CBS, A*, WHCA*. Focus on the trade-off between optimality
% and computational cost that motivates prioritised planning.

\todo{Cover CBS and variants briefly. Explain why optimal MAPF
solvers are impractical for real-time warehouse settings.}


% ── 2.2.2 ──────────────────────────────────────────────────────
\subsection{Lifelong and Warehouse MAPF}
\label{subsec:bg_mapf_lifelong}

% Lifelong MAPF: tasks arrive over time, agents never stop.
% This is the setting your system operates in.
% Key reference: Lifelong Multi-Agent Path Finding in Large-Scale
% Warehouses (Ma et al. / RHCR paper).

\todo{Describe lifelong MAPF~\citep{li2021lifelong}.
Discuss throughput as the relevant metric (not makespan per episode).
Note how your problem differs: task allocation is also decentralised,
not just path planning.}


% ── 2.2.3 ──────────────────────────────────────────────────────
\subsection{Prioritised Planning and Reservation Tables}
\label{subsec:bg_prioritised}

% Prioritised A* with space-time reservation tables: the approach
% you actually implement. Cover: completeness trade-offs,
% priority ordering, replanning triggers.

\todo{Describe prioritised planning and time-based reservation
tables. Explain the completeness caveat (not guaranteed with
fixed priorities). Justify why it is appropriate for your
real-time setting.}


% ── 2.3 ─────────────────────────────────────────────────────────
\section{Integrated Task Allocation and Path Planning}
\label{sec:bg_integrated}

% This section is where you identify the gap your thesis fills.
% Most work treats allocation and path planning as separate problems.
% Work that integrates them (DNA, DECNTR, Omni) is more relevant.

\todo{Review papers that integrate task allocation with path
planning: DNA~\citep{dna}, DECNTR~\citep{decntr},
Omni~\citep{omni}. Identify the common gap: most assume static
communication graphs or fully connected networks. Your work
explicitly addresses dynamic, potentially disconnected graphs.}


% ── 2.4 ─────────────────────────────────────────────────────────
\section{Communication-Constrained Coordination}
\label{sec:bg_communication}

% Cover work that explicitly models limited communication.
% Key point: most CBBA/GCBBA literature assumes the graph is given
% and fixed. Dynamic communication graphs (proximity-based) are
% less studied.

\todo{Review literature on coordination under communication
constraints. Discuss proximity-based communication graphs and
their impact on algorithm performance. Identify what makes
your dynamic replanning approach novel in this context.}


% ── 2.5 ─────────────────────────────────────────────────────────
\section{Warehouse Automation Systems}
\label{sec:bg_warehouse}

% Briefly ground the work in real warehouse systems (Kiva/Amazon,
% Geek+, etc.) to establish practical relevance without over-claiming
% industrial applicability.

\todo{Brief section on warehouse automation: induct/eject station
architecture, fleet management, why coordination failures
(deadlocks, incomplete allocation) are costly. Keep to 1 page max.}


% ── 2.6 ─────────────────────────────────────────────────────────
\section{Positioning of This Work}
\label{sec:bg_positioning}

% A focused 1--2 paragraph synthesis that maps the gap in the
% literature to exactly what your thesis contributes.
% Refer back to your RQs from Chapter 1.

\todo{Write a crisp positioning paragraph. The key claim:
existing work either (a) integrates allocation and path planning
but assumes full connectivity, or (b) handles disconnected graphs
but only for static allocation. Your thesis combines both:
dynamic GCBBA with event-driven replanning under arbitrary
communication topology.}
